\section{Conclusion}

The emergence of self-evolving agents marks a paradigm shift in artificial intelligence, moving beyond static, monolithic models toward dynamic agentic systems capable of continual learning and adaptation. As language agents are increasingly deployed in open-ended, interactive environments, the ability to evolve, adapting reasoning processes, tools, and behaviors in response to new tasks, knowledge, and feedback, has become essential for building the next generation of agentic systems. In this survey, we provide the first comprehensive and systematic review of self-evolving agents, organized around three foundational questions: \textit{what aspects of an agent should evolve, when evolution should occur, and how to implement evolutionary processes effectively}. Moreover, we discuss several methods for evaluating the progress of self-evolving agents in terms of metrics and benchmarks, followed by corresponding applications and future directions. 
The evolution of these agents will require significant advancements in models, data, algorithms, and evaluation practices, and so on.  Addressing issues such as catastrophic forgetting, human preference alignment during autonomous evolution, and the co-evolution of agents and environments will be key to unlocking agents that are not only adaptive but also trustworthy and aligned with human values. We hope this survey provides a foundational framework for researchers and practitioners to design, analyze, and advance the development and progress of self-evolving agents.

